%==============================================================================
\chapter{Description fonctionnelle des besoins}
%==============================================================================

\section{Gestion des entités (CRUD)}
Le système doit assurer la création, la consultation, la modification et la
suppression de toutes les entités métier : fermiers, fermes, sites, ruches,
corps/hausses, cadres, agents, visites et rapports, dans le respect des règles
de gestion énoncées au chapitre~3.

\section{Planification et suivi des visites}
\begin{itemize}
  \item Les visites sont \textbf{régulières ou irrégulières}, pour des motifs variés :
        inspection sanitaire, peuplement par un nouvel essaim, changement de reine,
        ravitaillement de nourriture, prescription de médicament, subdivision d'essaim,
        ajout de cadre et/ou d'étage d'extension, collecte de cadres de miel,
        évaluation de la production, évaluation de l'effectif et du volume de l'essaim.
  \item Chaque agent suit un \textbf{planning de visite approuvé} par l'agent superviseur de la ruche concernée.
\end{itemize}

\subsection{Rapport de visite}
Chaque visite donne lieu à un rapport contenant : \textbf{date, heure, durée,
raison, constatations, actions prévues, actions effectuées, recommandations},
une évaluation qualitative de l'effectif / du volume de l'essaim, son état de
santé et son niveau de productivité sur une \textbf{échelle de 1 à 3}.

\section{Tableaux de bord}

\subsection{Calendrier / matrice agents × ruches}
Matrice croisant \textbf{agents et ruches}, déclinée par mois, semaine, saison
et année, avec \textbf{filtres} et regroupements par site, par agent et par agent
responsable. Un clic sur une case représentant une visite planifiée affiche un
\textbf{pop-up} résumant ce qui est prévu (date, heure, raison, actions) et, si
la visite est déjà exécutée, ce qui a été réalisé ainsi que la date et l'heure
effectives.

\subsection{Tableau de bord production (fermier / propriétaire)}
Met en évidence les couples \textbf{site × ruche} dont la production (poids) a
dépassé un \textbf{seuil paramétrable}, signalant soit une collecte de miel
imminente, soit une intervention d'extension et/ou de division d'essaim pour
favoriser l'expansion de la colonie.

\subsection{Tableau de bord alertes (responsable)}
Signale les couples \textbf{site × ruche} dont la production ou l'effectif
baissent de façon inhabituelle au-delà d'un seuil paramétrable, déclenchant une
\textbf{alerte} et une inspection imminente.

\subsection{Tableau de bord de synthèse (propriétaire)}
Graphiques, courbes et histogrammes représentant :
\begin{itemize}
  \item la quantité produite par ferme et par site
  \item le nombre d'interventions par ferme, par site et par ruche
  \item le \textbf{retour sur investissement} (rapport coûts / production) pour arbitrer
        le sort de chaque site ou ruche (clôture, déplacement).
\end{itemize}

\section{Indicateurs de performance (volet automatisé : préparation)}
Les mesures sont \textbf{horodatées et géolocalisées}, exprimées dans des unités
\textbf{paramétrables} (internationalisation des systèmes d'unités). Les mesures
d'un même indicateur ne partagent pas nécessairement la même unité, les capteurs
étant hétérogènes. Indicateurs concernés :
\begin{itemize}
  \item poids du cadre, du socle et de l'étage d'extension
  \item température et humidité, à l'intérieur et à l'extérieur de la ruche
  \item effectif des mouvements d'entrée et de sortie des abeilles
  \item niveau de bruit / bourdonnement, intérieur et extérieur
  \item niveau de luminosité, intérieur et extérieur
  \item vitesse et direction du vent à l'extérieur
  \item état d'ouverture de la ruche (intervention, orage, vol, attaque animale, etc.).
\end{itemize}

\subsection{Analyse prédictive et surveillance connectée}
Les solutions de surveillance connectée (precision beekeeping) telles qu'Arnia,
BroodMinder ou BeeHero montrent que l'exploitation avancée des mêmes indicateurs
permet d'anticiper les événements de la colonie. Le modèle de données doit rester
ouvert à ces traitements, prévus pour le volet automatisé.
\begin{itemize}
  \item détection acoustique de l'essaimage 1 à 5 jours à l'avance par analyse de la signature sonore de la colonie
  \item confirmation de la présence de la reine à partir du bourdonnement
  \item détection d'ouverture, de choc ou de basculement de la ruche par accéléromètre (vol, chute, attaque animale), en complément de l'indicateur d'ouverture
  \item orientation de la ruche par boussole électronique
  \item pesée continue par balance pour suivre le flux de nectar et la dynamique de production
  \item analyse dans le cloud par intelligence artificielle des signatures acoustiques pour corréler son et activité de la ruche.
\end{itemize}

\subsection{Protocole d'ingestion des mesures}
Le volet automatisé reste hors périmètre de développement, mais l'interface
d'ingestion est spécifiée dès maintenant pour que le modèle de données et
l'architecture l'anticipent.
\begin{itemize}
  \item \textbf{Transport} : une passerelle terrain transmet les mesures au
        serveur, au choix selon la connectivité, par \textbf{REST sur HTTPS}
        (\texttt{POST /api/mesures}) ou par \textbf{MQTT} (publication sur le
        topic \texttt{zumm/\{rucheId\}/\{indicateur\}}).
  \item \textbf{Format} : \textbf{JSON}, une mesure portant l'identifiant de la
        ruche, le type d'indicateur, la valeur, l'unité, l'horodatage et la
        géolocalisation ; l'envoi par lot est un tableau de mesures.
  \item \textbf{Fréquence} : paramétrable par indicateur (par exemple poids
        toutes les 15 min, température toutes les 5 min, acoustique échantillonnée
        en continu), afin de maîtriser la volumétrie.
  \item \textbf{Robustesse} : ingestion \textbf{asynchrone} par file d'attente ;
        \textbf{idempotence} sur la clé (ruche, indicateur, horodatage) pour
        rejouer sans doublon ; en zone reculée, mise en tampon locale sur la
        passerelle puis synchronisation différée.
  \item \textbf{Sécurité} : passerelle authentifiée (clé ou certificat), canal
        TLS ; rejet des valeurs hors plage physique plausible (garde-fou).
\end{itemize}
\begin{verbatim}
POST /api/mesures        Content-Type: application/json
{
  "rucheId": 42, "typeIndicateur": "TEMPERATURE_INTERNE",
  "valeur": 34.8, "unite": "degC",
  "horodatage": "2026-07-09T14:05:00Z",
  "latitude": 36.806, "longitude": 10.181
}
\end{verbatim}

\subsection{Seuils par défaut et logique d'alerte}
Les seuils sont stockés dans la table \texttt{SEUIL} (paramétrables, chapitre~5).
Les valeurs ci-dessous sont des \textbf{valeurs par défaut de référence}, ancrées
dans le protocole apicole (chapitre~12).
\begin{longtable}{>{\bfseries}p{4.4cm} p{5.2cm} p{4.4cm}}
\toprule
\textbf{Indicateur} & \textbf{Seuil par défaut} & \textbf{Alerte déclenchée} \\
\midrule
\endhead
Température interne (couvain) & \textless\,32\,°C ou \textgreater\,36\,°C & Risque thermique sur le couvain \\
Humidité interne & \textgreater\,80\,\% & Excès d'humidité, risque sanitaire \\
Poids : gain de la hausse & \textgreater\ seuil de production & Collecte de miel imminente \\
Poids : chute soudaine & \textgreater\,1,5\,kg en moins d'1\,h & Essaimage, vol, chute ou attaque \\
Production ou effectif & Baisse relative \textgreater\,20\,\% vs période précédente & Inspection sanitaire imminente \\
Activité entrée/sortie & $\approx$\,0 en journée favorable & Colonie affaiblie ou reine absente \\
État d'ouverture & Ouvert hors visite planifiée & Intrusion ou perturbation \\
Signature acoustique & Motif d'essaimage détecté & Pré-alerte d'essaimage (1 à 5 j) \\
\bottomrule
\end{longtable}
\begin{encadre}
\textbf{Règle d'évaluation :} à chaque mesure, la valeur est convertie vers
l'unité de référence (formule de conversion, §5.3), puis comparée au seuil
paramétrable. Pour limiter les \emph{fausses alertes}, une alerte n'est levée que
si le dépassement \textbf{persiste sur $N$ mesures consécutives} (anti-rebond /
hystérésis). Toute alerte reste une \textbf{aide à la décision} confirmée par le
rapport de visite : l'agent conserve la validation finale.
\end{encadre}

\section{Grille de synthèse des fonctions}
\begin{longtable}{>{\bfseries}p{3.4cm} p{7.6cm} C{2.2cm}}
\toprule
\textbf{Fonction} & \textbf{Description} & \textbf{Priorité} \\
\midrule
\endhead
CRUD entités        & Gestion complète du modèle métier                       & Haute \\
Planning de visites & Planification + approbation par le superviseur           & Haute \\
Rapports de visite  & Saisie structurée des constats et actions               & Haute \\
Calendrier agents×ruches & Matrice filtrable avec pop-up de résumé            & Haute \\
Alertes production  & Détection des dépassements de seuil                     & Haute \\
Alertes sanitaires  & Détection des baisses anormales d'effectif/production   & Haute \\
Synthèse \& ROI     & Graphiques coûts/production, aide à la décision         & Moyenne \\
Indicateurs capteurs & Modèle de données prêt pour l'automatisation            & Moyenne \\
\midrule
\multicolumn{3}{l}{\textit{\color{mielfonce} Exigences anti-défaut (tirées des limites des solutions existantes)}} \\
\midrule
Déploiement autonome & Installation J2EE sans abonnement ni service tiers obligatoire & Haute \\
Mode hors-ligne & Saisie sans connexion et synchronisation différée & Haute \\
Authentification simple & OAuth Google et parcours de saisie terrain rapide & Haute \\
Modularité configurable & Activation des modules par profil via ConfigZumm.ini & Moyenne \\
Portabilité des données & Export CSV et TXT, service web API et sauvegarde, sans verrouillage & Haute \\
Aide contextuelle & Interface simple et cohérente, guidage de l'utilisateur & Moyenne \\
Parité web et mobile & Mêmes fonctions sur tous les supports & Moyenne \\
Internationalisation & Fichier XML FR, EN et AR, extensible à d'autres langues & Haute \\
Ouverture matérielle & Capteurs hétérogènes et unités paramétrables & Moyenne \\
Ingestion asynchrone & Mesures horodatées différées, volet manuel autonome & Moyenne \\
Validation humaine & Seuils paramétrables et confirmation par le rapport de visite & Haute \\
Sécurité des échanges & SSL et certificats X.509, accès authentifiés & Haute \\
Aide à la décision & Alertes indicatives, l'agent conserve la validation finale & Moyenne \\
\bottomrule
\end{longtable}

\section{Fonctionnalités complémentaires inspirées des solutions du marché}
Afin de renforcer l'ergonomie et la valeur d'usage, le système s'inspire des
bonnes pratiques des applications apicoles de référence telles que HiveTracks.
Ces fonctionnalités étendent les besoins déjà décrits sans les remplacer.
\begin{longtable}{>{\bfseries}p{4.2cm} p{7.5cm} C{2.1cm}}
\toprule
\textbf{Fonction} & \textbf{Description} & \textbf{Priorité} \\
\midrule
\endhead
Photos d'inspection & Attacher une ou plusieurs photos à chaque rapport de visite pour un suivi visuel de la colonie et des cadres. & Haute \\
Saisie vocale assistée & Dicter le compte rendu de visite, la transcription alimentant automatiquement les champs du rapport. & Basse \\
Contexte météo & Afficher la météo locale de la semaine à partir de la localisation du site pour éclairer les décisions d'intervention. & Moyenne \\
Carte et rayons de butinage & Positionner les ruches sur une carte et tracer des rayons de butinage (par exemple 1, 2 et 3 km, unités paramétrables) pour visualiser l'environnement accessible aux abeilles. & Moyenne \\
Liste de tâches et rappels & Gérer une liste de tâches et un calendrier de rappels (nourrissement, inspection, statut de la reine) pour ne manquer aucune échéance. & Haute \\
Suivi de la reine & Enregistrer l'historique du statut de la reine (présence, âge, remplacement, marquage) au fil des visites. & Haute \\
Récolte et QR code & Saisir les récoltes de miel et générer un QR code par lot présentant les notes de dégustation, la provenance et les photos, pour la traçabilité et la valorisation. & Moyenne \\
Accès multi-supports & Offrir un accès web et une expérience adaptée aux terminaux mobiles pour la saisie terrain. & Moyenne \\
\bottomrule
\end{longtable}
\begin{encadre}
\textbf{Cohérence avec les exigences :} les rayons de butinage et la météo
réutilisent la géolocalisation des sites et le mécanisme d'unités paramétrables.
Les rappels s'appuient sur le planning de visite existant, et le suivi de la
reine enrichit le rapport de visite déjà défini.
\end{encadre}

\section{Fonctions de gestion avancée (inspiration Apiary Book et BeePlus)}
Les applications de gestion apicole Apiary Book et BeePlus mettent en avant des
fonctions d'organisation et de traçabilité qui complètent le périmètre du
système.
\begin{longtable}{>{\bfseries}p{4.2cm} p{7.5cm} C{2.1cm}}
\toprule
\textbf{Fonction} & \textbf{Description} & \textbf{Priorité} \\
\midrule
\endhead
Mode hors-ligne et synchronisation & Saisir les données terrain sans connexion et les synchroniser automatiquement au retour du réseau. & Haute \\
Suivi des traitements & Conserver l'historique des médicaments et traitements sanitaires appliqués à chaque ruche. & Haute \\
Gestion du matériel et inventaire & Suivre le matériel (corps, hausses, cadres, ruches) et son affectation aux sites. & Moyenne \\
Suivi financier & Enregistrer coûts et revenus par ferme, site et ruche, en appui du calcul du retour sur investissement. & Moyenne \\
Export des données & Exporter les enregistrements aux formats CSV et TXT pour une analyse externe. & Moyenne \\
Détection de maladies & Aider à l'identification des maladies à partir des constatations de visite. & Basse \\
Sauvegarde des données & Sauvegarder et restaurer les données dans le cloud. & Moyenne \\
\bottomrule
\end{longtable}
\begin{encadre}
\textbf{Cohérence avec les exigences :} le suivi financier alimente le tableau de
bord de retour sur investissement, l'export CSV prolonge le service web API, et
le suivi des traitements enrichit le motif de visite prescription de médicament.
\end{encadre}

\section{Limites connues des solutions existantes et exigences pour les éviter}
L'analyse des retours d'expérience sur les solutions du marché (HiveTracks,
Apiary Book, BeePlus) et sur les systèmes de surveillance connectée (Arnia,
BroodMinder, BeeHero) fait ressortir des limites récurrentes. Le tableau suivant
met chaque limite en regard de l'exigence retenue pour l'éviter dans Zümm.
\begin{longtable}{>{\bfseries}p{5.2cm} p{9.3cm}}
\toprule
\textbf{Limite observée} & \textbf{Exigence retenue pour l'éviter} \\
\midrule
\endhead
Abonnement payant et paywall & Système autodéployable en J2EE, sans abonnement obligatoire, l'ensemble des fonctions de gestion reste accessible. \\
Fonctionnement hors-ligne limité & Mode hors-ligne robuste avec synchronisation différée, la saisie terrain n'exige pas de connexion permanente. \\
Création de compte imposée et friction & Authentification simple par OAuth Google et parcours de saisie rapide, l'ergonomie prime. \\
Surcharge de fonctions inutiles & Modules activables par profil et par le fichier de configuration ConfigZumm.ini, grâce à l'architecture faiblement couplée. \\
Dépendance et perte des données & Données maîtrisées par l'utilisateur, export CSV et TXT, service web API et sauvegarde, sans verrouillage propriétaire. \\
Interface dense et courbe d'apprentissage & Interface simple, cohérente et conviviale avec aide contextuelle, conformément aux exigences d'ergonomie. \\
Fonctions inégales selon la plateforme & Parité fonctionnelle entre accès web et mobile grâce à l'architecture multi-niveaux. \\
Solutions uniquement anglophones & Internationalisation par fichier XML, au moins FR, EN et AR, ouverte à d'autres langues. \\
Coût élevé du matériel de capteurs & Ouverture aux capteurs hétérogènes et bon marché avec unités paramétrables, sans verrouillage matériel. \\
Autonomie et connectivité en zone reculée & Ingestion asynchrone de mesures horodatées, le volet manuel fonctionne sans capteurs. \\
Fausses alertes et fiabilité des capteurs & Seuils paramétrables et validation humaine par le rapport de visite, les alertes restent une aide à la décision. \\
Confidentialité et sécurité des données & Chiffrement SSL et certificats X.509, authentification des accès. \\
Surconfiance dans l'automatisation & Positionnement en aide à la décision, l'agent apiculteur conserve la validation finale. \\
\bottomrule
\end{longtable}

